\documentclass[11pt,a4paper]{article}

% Preâmbulo e estilos compactos (garante limite estrito de 4 páginas de conteúdo)
% abntformat.tex
% ------------------------------------------------------------------------------
% Configuração de Pacotes e Estilo Visual para Relatório Técnico (Máx. 4 páginas)
% ------------------------------------------------------------------------------
\usepackage[utf8]{inputenc}
\usepackage[T1]{fontenc}
\usepackage[brazil]{babel}
\usepackage{lmodern}

% Margens de 2.0cm para aproveitamento ideal do espaço no limite de 4 páginas
\usepackage{geometry}
\geometry{
    a4paper,
    margin=2.0cm,
    top=2.0cm,
    bottom=2.0cm
}

\usepackage{graphicx}
\usepackage{subcaption}
\usepackage{booktabs}
\usepackage{siunitx}
\usepackage{amsmath, amssymb}
\usepackage{hyperref}
\usepackage{enumitem}
\usepackage{titlesec}
\usepackage{xcolor}

% Configuração de pontuação numérica
\sisetup{
    output-decimal-marker = {,},
    per-mode = symbol
}

% Espaçamento compacto em listas para não inflar páginas desnecessariamente
\setlist{nosep, leftmargin=1.5em, topsep=2pt, partopsep=0pt}

% Ajuste sutil de espaçamento de seções para garantir o teto de 4 páginas
\titlespacing*{\section}{0pt}{1.2ex plus 0.4ex minus 0.2ex}{0.6ex plus 0.2ex}
\titlespacing*{\subsection}{0pt}{1.0ex plus 0.3ex minus 0.2ex}{0.4ex plus 0.2ex}

% Configuração de hiperlinks
\hypersetup{
    colorlinks=true,
    linkcolor=black,
    citecolor=black,
    urlcolor=blue
}


\begin{document}

% ==============================================================================
% Capa Institucional ABNT / UFPR (página separada não contabilizada nas 4 páginas)
% ==============================================================================
% capa.tex
% ==============================================================================
% CAPA ESTILO ABNT / UFPR - CI1061 REDES DE COMPUTADORES II
% ==============================================================================
\begin{titlepage}
    \thispagestyle{empty}

    \begin{center}
        % Topo: instituição, setor e departamento
        {\bfseries UNIVERSIDADE FEDERAL DO PARANÁ\par}
        {\bfseries SETOR DE CIÊNCIAS EXATAS\par}
        {\bfseries DEPARTAMENTO DE INFORMÁTICA\par}
        {\bfseries BACHARELADO EM CIÊNCIA DA COMPUTAÇÃO\par}

        \vspace*{2.5cm}

        % Nomes dos integrantes da equipe
        {\bfseries ANDRÉ AKIRA ...\par}
        GRR2024... -- \texttt{login}\par
        \vspace{0.4cm}
        {\bfseries LAISA MARCELINO SANTOS RODRIGUES\par}
        GRR20243783 -- \texttt{lmsr24}\par
        \vspace{0.4cm}
        {\bfseries MURILO ...\par}
        GRR2024... -- \texttt{login}\par

        \vspace*{4.0cm}

        % Título do trabalho
        {\bfseries\Large
        RELATÓRIO TÉCNICO: SISTEMA DE BALANCEAMENTO DE CARGAS TCP (LOAD BALANCER)\par}
        \vspace{0.5cm}
        {\large Disciplina: Redes de Computadores II -- CI1061\par}
        {\large Professor: Dr. Giovanni Venâncio\par}

        \vspace*{\fill}

        Curitiba -- PR\par
        2026
    \end{center}
\end{titlepage}

\clearpage

% Reinicia a numeração arábica a partir da primeira página de texto efetivo
\pagenumbering{arabic}
\setcounter{page}{1}

% ==============================================================================
\section{Introdução e Visão Geral da Arquitetura}
% ==============================================================================

Em sistemas distribuídos modernos, o balanceamento de carga na camada de transporte é fundamental para prover escalabilidade horizontal e alta disponibilidade. A concentração de todas as conexões em uma única máquina servidora gera um gargalo computacional e cria um ponto único de falha (\textit{Single Point of Failure} -- SPOF). O presente trabalho prático desenvolve um sistema completo de balanceamento de carga TCP concorrente para a disciplina de Redes de Computadores II (CI1061).

O sistema é composto por cinco componentes essenciais, operando exclusivamente com sockets TCP padrão da biblioteca nativa do Python:
\begin{enumerate}
    \item \textbf{Clientes Concorrentes:} Simulam $M$ conexões paralelas que enviam fluxos de mensagens terminadas em \texttt{"\textbackslash n"} durante $K$ segundos. Em estrita conformidade com o princípio de transparência de rede, os clientes conhecem \textit{apenas} o endereço IP e a porta do Load Balancer, sendo terminantemente proibida a comunicação direta com os servidores backend.
    \item \textbf{Load Balancer (LB):} Agente intermediário concorrente e não bloqueante que escuta conexões de entrada, consulta a tabela de servidores, despacha cada fluxo para um nó backend conforme a política configurada e opera como proxy TCP bidirecional (\textit{relay}) full-duplex.
    \item \textbf{Servidores Backend ($N$ instâncias):} Servidores TCP concorrentes independentes que recebem o tráfego repassado pelo LB, processam a requisição e devolvem a confirmação \texttt{"OK\textbackslash n"}. Estão sujeitos a panes forçadas e recuperações dinâmicas.
    \item \textbf{Tabela de Servidores:} Estrutura de dados compartilhada em memória que mantém o número total global de conexões mantidas pelo LB, a contagem de conexões ativas e acumuladas de cada servidor, endereço IP, porta e status operacional (\texttt{UP} ou \texttt{DOWN}).
    \item \textbf{Sistema de Monitoramento (\textit{Health Check}):} Mecanismo periódico executado em thread dedicada que diagnostica a integridade dos servidores, atualizando a tabela diante de quedas e recuperações.
\end{enumerate}

\begin{figure}[!ht]
\centering
\small
\begin{verbatim}
+---------------+          +---------------------------------+          +---------------+
|   Clientes    |   TCP    |          LOAD BALANCER          |   TCP    | Servidor 1..N |
|  (1..M, K s)  | =======> | - Proxy TCP (Relay Bidirecional)| =======> | (portas 800N) |
+---------------+  Dados   | - Políticas (Round-Robin / LC)  |  Dados   | Resposta "OK" |
                           | - Tabela de Servidores (Mutex)  |          +---------------+
                           | - HealthMonitor (Sondagens TCP) | - - - - > Controle: UP /
                           +---------------------------------+  Sonda    DOWN periódico
\end{verbatim}
\vspace{-0.6em}
\caption{Arquitetura do sistema: canais de dados (linhas duplas) e controle de saúde (tracejado).}
\label{fig:arquitetura}
\end{figure}

% ==============================================================================
\section{Decisões de Projeto e Detalhes de Implementação}
% ==============================================================================

A implementação foi construída em Python 3 utilizando unicamente os módulos padrão \texttt{socket}, \texttt{threading}, \texttt{argparse}, \texttt{time} e \texttt{signal}.

\subsection{Modelo de Concorrência: Threads vs. Processos}
Tanto no Load Balancer quanto nos servidores backend, a carga de trabalho é predominantemente limitada por entrada e saída de rede (\textit{I/O-bound}). Por essa razão, optou-se pelo modelo \textit{thread-per-connection}. Ao receber uma nova conexão via chamada bloqueante \texttt{accept()}, o socket despacha imediatamente o descritor para uma thread trabalhadora criada em modo \textit{daemon}.

Essa estratégia proporciona baixo consumo de memória e chaveamento de contexto ágil no Linux, além de viabilizar o compartilhamento natural da tabela de servidores em memória entre o proxy e o monitor de saúde. No CPython, operações bloqueantes de rede (\texttt{socket.recv()} e \texttt{socket.sendall()}) liberam explicitamente o \textit{Global Interpreter Lock} (GIL) na chamada de sistema do kernel, permitindo paralelismo de I/O entre múltiplos clientes.

\subsection{Tabela de Servidores e Sincronização Atômica}
A classe \texttt{ServerTable} centraliza o estado do cluster. Cada nó é modelado por um objeto \texttt{ServerEntry}, contendo os atributos \texttt{id}, \texttt{ip}, \texttt{port}, \texttt{status}, \texttt{active\_connections} e \texttt{total\_connections}.

O acesso concorrente simultâneo por dezenas de threads de clientes e pela thread de monitoramento impõe riscos de condições de corrida (\textit{race conditions}). Para assegurar integridade, todas as leituras e mutações são protegidas por uma primitiva de exclusão mútua (\texttt{threading.Lock}).

Em particular, identificou-se uma janela de vulnerabilidade entre o instante em que o LB consulta o servidor com menor carga e o momento em que a conexão é efetivamente iniciada. Para sanar esse problema, desenvolveu-se o método atômico \texttt{select\_and\_reserve(policy)}: sob a proteção do mesmo lock, o método avalia a política, escolhe o servidor disponível e já incrementa imediatamente tanto suas conexões ativas quanto o contador global do LB, eliminando o risco de sobrecarga concentrada durante rajadas de conexões.

\subsection{Proxy TCP e Relay Bidirecional Full-Duplex}
O repasse de tráfego ocorre na função \texttt{relay(src, dst)}. Para garantir transporte bidirecional concorrente e não bloqueante, cada conexão cliente-LB-backend dispara duas threads de encaminhamento:
\begin{itemize}
    \item \textbf{Canal de Entrada:} lê bytes do cliente e transmite integralmente ao servidor backend;
    \item \textbf{Canal de Retorno:} lê a confirmação do servidor e transmite diretamente ao cliente.
\end{itemize}
Quando um dos lados encerra o envio, invoca-se \texttt{dst.shutdown(socket.SHUT\_WR)}, drenando ordenadamente os dados residuais. Ao término do fluxo, o bloco \texttt{finally} assegura o fechamento seguro de ambos os sockets e invoca \texttt{table.decrement\_connection(backend.id)}, garantindo o estorno da conexão ativa mesmo em casos de desconexão abrupta.

\subsection{Monitoramento Contínuo e Tolerância a Falhas}
A classe \texttt{HealthMonitor} (em \texttt{src/monitor.py}) executa em uma thread autônoma em segundo plano com intervalo configurável (padrão de \SI{2.0}{\second}). A sondagem de integridade é implementada via tentativa de conexão TCP rápida (\texttt{socket.connect}) com timeout curto de \SI{1.0}{\second}:
\begin{itemize}
    \item \textbf{Sucesso na conexão:} o socket de teste é fechado de imediato e o servidor é marcado como \texttt{UP}. Caso estivesse \texttt{DOWN}, registra-se o evento de recuperação com timestamp nos logs.
    \item \textbf{Falha (\texttt{ConnectionRefused} ou timeout):} o servidor é marcado atomicamente como \texttt{DOWN}, sendo imediatamente excluído das decisões de roteamento do LB.
\end{itemize}
O projeto implementa tratamento de falhas em duas vias: manual, por meio do script \texttt{scripts/recover\_server.sh} (que reinicia a instância na mesma porta via verificação de porta livre), e automática, através da opção \texttt{--auto-recover} do monitor, que relança o processo como subprocesso do sistema com proteção contra relançamentos em loop.

% ==============================================================================
\section{Implementação das Políticas de Balanceamento}
% ==============================================================================



% ==============================================================================
\section{Metodologia e Resultados Experimentais}
% ==============================================================================



% ==============================================================================
\section{Conclusão}
% ==============================================================================


% refs.tex
% ==============================================================================
% REFERÊNCIAS BIBLIOGRÁFICAS
% ==============================================================================
\section*{Referências}
\addcontentsline{toc}{section}{Referências}

\begin{itemize}[leftmargin=*]
    \item KUROSE, James F.; ROSS, Keith W. \textit{Redes de Computadores e a Internet: Uma Abordagem Top-Down}. 8. ed. São Paulo: Pearson Education do Brasil, 2021.

    \item TANENBAUM, Andrew S.; WETHERALL, David J. \textit{Redes de Computadores}. 5. ed. São Paulo: Pearson Prentice Hall, 2011.

    \item STEVENS, W. Richard. \textit{TCP/IP Illustrated, Volume 1: The Protocols}. Reading: Addison-Wesley, 1994.

    \item STEVENS, W. Richard; FENNER, Bill; RUDOFF, Andrew M. \textit{UNIX Network Programming, Volume 1: The Sockets Networking API}. 3. ed. Boston: Addison-Wesley, 2004.

    \item PYTHON SOFTWARE FOUNDATION. \textit{socket --- Low-level networking interface}. Python 3 Standard Library Documentation. Disponível em: \url{https://docs.python.org/3/library/socket.html}. Acesso em: nov. 2026.

    \item VENÂNCIO, Giovanni. \textit{Especificação do Trabalho Prático: Sistema de Balanceamento de Cargas}. Disciplina de Redes de Computadores II (CI1061), Departamento de Informática, Universidade Federal do Paraná (UFPR), 2026/2.
\end{itemize}


\end{document}
